% Figuras vetoriais autorais: parâmetros controlam etapas da apresentação.
\definecolor{eletron}{RGB}{0,83,134}
\definecolor{lacuna}{RGB}{199,123,0}
\definecolor{ativo}{RGB}{0,132,75}
\tikzset{fio/.style={line width=.9pt,black!65},ligado/.style={line width=2pt,ativo},cortado/.style={line width=.7pt,black!25},anotacao/.style={font=\small,align=center},passo/.style={font=\scriptsize,text=blue_theme}}
\newcommand{\chamada}[1]{\par\medskip\centering\textcolor{blue_theme}{\textbf{#1}}\par}
\newcommand{\etapa}[2]{\par\smallskip{\scriptsize\color{blue_theme}Etapa #1\enspace•\enspace #2}\par}
% PN: 0 contato inicial; 1 difusão; 2 equilíbrio; 3 direta; 4 reversa.
\newcommand{\juncaoPN}[1]{%
\begin{tikzpicture}[x=.87cm,y=.85cm,font=\small]
\path[use as bounding box] (-.3,-1.65) rectangle (10.3,4.9);
\fill[lacuna!9] (0,0) rectangle (5,3);\fill[eletron!9] (5,0) rectangle (10,3);
\node[blue_theme] at (2,3.5){\textbf{Região P}};\node[blue_theme] at (8,3.5){\textbf{Região N}};
\ifnum#1>1
 \def\depl{1.25}\ifnum#1=3\def\depl{.6}\fi\ifnum#1=4\def\depl{2}\fi
 \fill[gray!25] (5-\depl,0) rectangle (5+\depl,3);
 \draw[densely dashed,gray] (5-\depl,0)--(5-\depl,3);\draw[densely dashed,gray] (5+\depl,0)--(5+\depl,3);
 \foreach \y in {.55,1.25,1.95,2.6}{\node[eletron] at (5-.5*\depl,\y){$-$};\node[lacuna] at (5+.5*\depl,\y){$+$};}
 \draw[<->,blue_theme] (5-\depl,4.15)--node[above]{\strut Depleção}(5+\depl,4.15);
 \draw[-{Latex},very thick,blue_theme] (5+\depl,-.55)--node[below]{Campo interno}(5-\depl,-.55);
\fi
\foreach \x/\y in {.7/.6,1.6/1.4,2.5/.65,.8/2.3,2.6/2.4}{\draw[lacuna,fill=white,line width=1pt] (\x,\y) circle(.13);}
\foreach \x/\y in {7.5/.6,8.4/1.4,9.3/.65,7.4/2.4,9.2/2.3}{\fill[eletron] (\x,\y) circle(.12);}
\ifnum#1<2
 \foreach \x/\y in {3.6/.8,4.3/2,3.4/2.5}{\draw[lacuna,fill=white,line width=1pt] (\x,\y) circle(.13);}
 \foreach \x/\y in {5.7/.8,6.4/2,6.6/2.5}{\fill[eletron] (\x,\y) circle(.12);}
\fi
\ifnum#1=1
 \draw[-{Latex},very thick,eletron] (6.3,1.2)--node[above]{elétrons}(3.8,1.2);
 \draw[-{Latex},very thick,lacuna] (3.6,2.3)--node[above]{lacunas}(6.1,2.3);
 \node[font=\scriptsize] at (5,-1.25){Difusão e recombinação próximas à interface};
\fi
\ifnum#1=0\node at (5,-1){\textcolor{lacuna}{$\circ$} lacuna \qquad\textcolor{eletron}{$\bullet$} elétron};\fi
\ifnum#1=2\node[font=\scriptsize] at (5,-1.3){Sinais na faixa cinza: íons fixos, não portadores móveis};\fi
\ifnum#1=3
 \node at (1.2,4.3){$+$};\node at (8.8,4.3){$-$};
 \node[font=\scriptsize] at (5,-1.3){Direta: barreira menor, corrente maior};
\fi
\ifnum#1=4
 \node at (1.2,4.3){$-$};\node at (8.8,4.3){$+$};
 \node[font=\scriptsize] at (5,-1.3){Reversa: barreira maior, pequena fuga antes da ruptura};
\fi
\end{tikzpicture}}
% NMOS planar: 0 sem canal forte; 1 depleção; 2 canal; 3 corrente.
\newcommand{\moscanal}[1]{%
\begin{tikzpicture}[x=.85cm,y=.8cm,font=\small]
\path[use as bounding box] (-.3,-1.1) rectangle (11,5.2);
\fill[eletron!8,draw=blue_theme] (0,0) rectangle (10,2.8);
\fill[yellow!35,draw=blue_theme] (.5,1.7) rectangle (2.5,2.8);
\fill[yellow!35,draw=blue_theme] (7.5,1.7) rectangle (9.5,2.8);
\node at (1.5,2.2){$n^+$};\node at (8.5,2.2){$n^+$};
\fill[gray!40,draw=blue_theme] (2.7,2.85) rectangle (7.3,3.08);
\fill[eletron!55,draw=blue_theme] (2.7,3.08) rectangle (7.3,3.55);
\node at (5,.5){Corpo P};\node at (1.5,3.9){Fonte};\node at (8.5,3.9){Dreno};\node at (5,4){Porta};
\draw (1.5,2.8)--(1.5,3.6);\draw (8.5,2.8)--(8.5,3.6);\draw (5,3.55)--(5,3.8);
\node[font=\scriptsize,anchor=west] at (8.1,4.72){Dielétrico};\draw[-{Latex},black!50] (8.1,4.35)--(7.65,4.35)--(7.25,2.97);
\ifnum#1<1
 \foreach \x/\y in {3.3/1.2,4.3/2.1,5.2/1.4,6.5/2.1}{\draw[lacuna,fill=white] (\x,\y) circle(.12);}
 \node at (5,4.75){$V_{GS}$ baixo};
\else
 \fill[gray!20] (2.5,1.55) rectangle (7.5,2.8);
 \foreach \x in {3,4,5,6,7}{\node[eletron,font=\scriptsize] at (\x,2.1){$-$};}
 \foreach \x in {3.2,4.4,5.6,6.8}{\node[red] at (\x,3.32){$+$};}
 \ifnum#1=1
  \node at (5,4.75){A tensão de porta aumenta};
  \draw[-{Latex},lacuna,thick] (4,1.55)--(4,.8);\draw[-{Latex},lacuna,thick] (6,1.55)--(6,.8);
  \node[font=\scriptsize] at (5,-.65){Lacunas se afastam da superfície: depleção};
 \else
  \fill[eletron!70] (2.5,2.57) rectangle (7.5,2.8);
  \foreach \x in {2.9,3.6,4.3,5,5.7,6.4,7.1}{\node[white,font=\scriptsize] at (\x,2.69){$-$};}
  \node at (5,4.75){$V_{GS}>V_{TN}$: canal de elétrons};
  \ifnum#1=3
   \draw[-{Latex},white,line width=3.5pt] (8.4,2.69)--(1.6,2.69);\draw[-{Latex},line width=2pt,ativo] (8.4,2.69)--(1.6,2.69);\node[ativo] at (5,1.15){Corrente convencional $I_D$};
   \node[font=\scriptsize] at (5,-.65){Com $V_D>V_S$, há corrente pelo canal};
  \else\node[font=\scriptsize] at (5,-.65){Inversão: a superfície conduz por elétrons};\fi
 \fi
\fi
\end{tikzpicture}}
% Circuitos por estado: caminho ativo verde; transistor bloqueado cinza.
\newcommand{\inversorestado}[1]{%
\begin{circuitikz}[scale=.9,transform shape]
\path[use as bounding box] (-3,-2.9) rectangle (3,3);
\ifnum#1=0\def\pstyle{ativo}\def\nstyle{black!25}\def\saida{1}\else\def\pstyle{black!25}\def\nstyle{ativo}\def\saida{0}\fi
\node[pmos,color=\pstyle] (p) at (0,1.3) {};\node[nmos,color=\nstyle] (n) at (0,-1.1) {};
\draw[ligado,color=\pstyle] (p.S)--(0,2.2) node[above,black]{$\VDD$};
\draw[ligado,color=\pstyle] (p.D)--(0,.1);
\draw[ligado,color=\nstyle] (n.D)--(0,.1);\draw[ligado,color=\nstyle] (n.S)--(0,-2.1) node[ground,black]{};
\draw[fio] (p.G)--(-1.5,1.3)--(-1.5,-1.1)--(n.G);\draw[fio] (-1.5,.1)--(-2.1,.1) node[left]{$A=#1$};
\draw[fio] (0,.1) to[short,*-] (1.3,.1) node[right]{$Y=\saida$};
\node[font=\small,text=\pstyle,right] at (.5,1.3){PMOS};\node[font=\small,text=\nstyle,right] at (.5,-1.1){NMOS};
\node[font=\scriptsize,ativo] at (0,-2.8){Verde: caminho que fixa a saída};
\end{circuitikz}}
\newcommand{\portaestado}[3]{%
\begin{circuitikz}[scale=.78,transform shape]
\path[use as bounding box] (-3,-3.8) rectangle (3.7,3.6);
\ifnum#2=0\def\pa{ativo}\def\na{black!25}\else\def\pa{black!25}\def\na{ativo}\fi
\ifnum#3=0\def\pb{ativo}\def\nb{black!25}\else\def\pb{black!25}\def\nb{ativo}\fi
\ifnum#1=0
 \pgfmathtruncatemacro{\out}{1-#2*#3}
 \node[pmos,color=\pa] (pa) at (-1.2,1.7) {};\node[pmos,color=\pb] (pb) at (1.4,1.7) {};
 \node[nmos,color=\na] (na) at (0,-.7) {};\node[nmos,color=\nb] (nb) at (0,-2.3) {};
 \draw[fio] (pa.S)--(-1.2,2.6)--(1.4,2.6)--(pb.S);\node[above] at (.1,2.6){$\VDD$};
 \draw[ligado,color=\pa] (pa.D)--(-1.2,.5)--(0,.5);\draw[ligado,color=\pb] (pb.D)--(1.4,.5)--(0,.5);
 \draw[ligado,color=\na] (0,.5)--(na.D);\draw[ligado,color=\na] (na.S)--(0,-1.5);
 \draw[ligado,color=\nb] (0,-1.5)--(nb.D);\draw[ligado,color=\nb] (nb.S)--(0,-3.2) node[ground,black]{};
 \draw[fio] (0,.5) to[short,*-] (2.2,.5) node[right]{$Y=\out$};
\else
 \pgfmathtruncatemacro{\out}{(1-#2)*(1-#3)}
 \node[pmos,color=\pa] (pa) at (0,2) {};\node[pmos,color=\pb] (pb) at (0,.4) {};
 \node[nmos,color=\na] (na) at (-1.2,-1.8) {};\node[nmos,color=\nb] (nb) at (1.4,-1.8) {};
 \draw[ligado,color=\pa] (pa.S)--(0,2.85) node[above,black]{$\VDD$};\draw[ligado,color=\pa] (pa.D)--(0,1.2);
 \draw[ligado,color=\pb] (0,1.2)--(pb.S);\draw[ligado,color=\pb] (pb.D)--(0,-.6);
 \draw[ligado,color=\na] (na.D)--(-1.2,-.6)--(0,-.6);\draw[ligado,color=\nb] (nb.D)--(1.4,-.6)--(0,-.6);
 \draw[ligado,color=\na] (na.S)--(-1.2,-2.8);\draw[ligado,color=\nb] (nb.S)--(1.4,-2.8);
 \draw[fio] (-1.2,-2.8)--(1.4,-2.8);\draw (0,-2.8) node[ground]{};
 \draw[fio] (0,-.6) to[short,*-] (2.2,-.6) node[right]{$Y=\out$};
\fi
\draw (pa.G) node[left,black]{$A=#2$};\draw (pb.G) node[left,black]{$B=#3$};
\draw (na.G) node[left,black]{$A=#2$};\draw (nb.G) node[left,black]{$B=#3$};
\end{circuitikz}}
